\chapter{Introduction}
\label{ch:introduction}


\section{Context and Motivation}

Sepsis is a life-threatening condition caused by a dysregulated host response
to infection \cite{singer2016sepsis3}. It accounts for approximately 11 million
deaths annually, nearly one fifth of deaths worldwide \cite{rudd2020global},
and each additional hour before antibiotic administration is associated with
higher mortality \cite{seymour2017time}. An early warning system that detects
sepsis before clinical recognition could allow treatment to start sooner.

This potential has motivated a large body of research, but high
discrimination has not translated reliably into clinical value. Wang et al.\
\cite{wang2025srpm} reviewed 91 published real-time sepsis prediction models
and found a median internal AUROC of 0.811. Among the models that reported it,
the median Utility Score, which rewards timely alerts and penalises late and
false ones, fell from $+0.381$ internally to $-0.164$ externally; a negative
score means that alerting does more harm than issuing no alerts at all. In
prospective deployment, Wong et al.\ \cite{wong2026epic} found encounter-level
AUROC between 0.82 and 0.92 for the Epic Sepsis Model v2 across 227,091
encounters at four US health systems, with positive predictive value of 0.13
to 0.26, or 21 to 35 alerts investigated per sepsis case identified. These
studies show that discrimination can remain high while external utility is
poor and alert burden substantial.

Study design is a further source of variation. Cohen et al.\
\cite{cohen2024subtle} showed on MIMIC-III that varying only the
operationalisation of Sepsis-3 changed the identified cohort from 867 to 2,178
patients, and that the label choice affected performance as much as the model
choice. Kamran et al.\ \cite{kamran2024evaluation} showed that published models
partly learn to predict when clinicians start treatment rather than when
patients deteriorate.

These findings shift the focus from model ranking to study design. This thesis
asks how much of the variation in reported performance is attributable to the
choice of model, and how much to four design choices: how sepsis is defined,
when prediction is made relative to treatment, how the data are split, and
which metric is reported. Each is treated as a separate source of variation
and its contribution is estimated. The pre-registered term \emph{variance
decomposition}, which the title reflects, refers here to a comparison of
one-factor-at-a-time performance ranges, not to a formal partition of
statistical variance (Section~\ref{sec:decomposition}).


\section{Limitations of the Conventional Approach}

The conventional approach to sepsis prediction selects one Sepsis-3
operationalisation, one train/test split and one headline metric, usually
AUROC, and then compares models. Most of the 91 models reviewed by Wang et al.\
\cite{wang2025srpm} follow this pattern, which has four limitations:

\begin{enumerate}
    \item \textbf{The sepsis label is not fixed.} Different operationalisations
    of Sepsis-3 applied to the same database produce substantially different
    cohorts and performance estimates \cite{cohen2024subtle}. The label choice
    can change results as much as the model choice.

    \item \textbf{The data split can produce optimistic estimates.} Guo et al.\
    \cite{guo2022temporal} found that sepsis prediction degraded more than any
    of the other clinical tasks they evaluated when tested on a later admission
    era. Random splits mix admission eras and may therefore produce optimistic
    estimates when clinical practice or documentation changes over time.

    \item \textbf{AUROC does not reflect clinical usefulness.} A model can
    retain a high AUROC while its clinical utility falls below that of issuing
    no alerts \cite{wang2025srpm}.

    \item \textbf{The label incorporates clinician behaviour.} Because
    Sepsis-3 is operationalised through culture and antibiotic orders, models
    may learn when clinicians suspect sepsis and begin treatment rather than
    when the patient's condition worsens \cite{kamran2024evaluation}. The time
    at which prediction is evaluated relative to treatment therefore matters.
\end{enumerate}

Comparing models while holding these design choices fixed leaves their
contributions unmeasured. This study varies them explicitly and estimates each
contribution with an uncertainty interval, to compare their effects on
reported performance with the effect of model choice.


\section{Research Question and Pre-Registered Hypotheses}
\label{sec:rq}

\noindent\textbf{Research question.} \textit{In real-time sepsis prediction on
MIMIC-IV, how much of the variation in reported performance is attributable to
(a) the choice of model, and (b) the researcher's decisions: how sepsis is
defined, how the data are split, when prediction is made, and how performance
is measured?}

All discrimination estimates concern detection of the recorded onset hour
among at-risk person-hours: the models estimate the hazard of onset in the
current hour, so they measure onset detection rather than prediction over a
future horizon. Because predictors and outcomes are binned hourly, the analysis
does not establish that measurements preceded onset within the scored hour
(Section~\ref{sec:framing}; Limitation~\ref{lim:current_hour}). Here,
``real-time'' denotes hourly dynamic estimation without retrospective
information, not demonstrated advance warning.

Six hypotheses were pre-registered before any of the modelling reported here:

\begin{description}
    \item[H1 (Label dominance)] Changing the sepsis definition affects AUROC at
    least as much as changing the model.
    \item[H2 (Coefficient instability)] At least one clinical feature changes its
    effect direction, or changes in magnitude by a factor of two or more, across
    label variants.
    \item[H3 (Treatment leakage)] Restricting evaluation to before treatment
    begins reduces AUROC by at least 0.03.
    \item[H4 (Split optimism)] Time-based testing gives at least 0.02 lower AUROC
    than random testing.
    \item[H5 (Metric divergence)] The Utility Score declines proportionally more
    than AUROC when moving from internal to external validation.
    \item[H6 (Model-class null)] The primary discrete-time hazard model is not
    outperformed by gradient-boosted trees by more than 0.02 AUROC, under a
    fixed label and temporal split.
\end{description}

H1 and H6 represent competing explanations of performance variation:
confirming H1 would make the label the main source, whereas rejecting H6 would
place more weight on the model.

Two hypotheses could not be evaluated in their pre-registered form. Under the
pre-registered onset rule the pre-treatment window of H3 contains no onsets by
construction (Section~\ref{sec:h3}), and the internal Utility Score is too
close to zero for the proportional decline of H5 to be defined, so an
internal-to-external AUROC contrast takes its place (Section~\ref{sec:h5}).
Section~\ref{sec:h6} sets out how the interval null of H6 is tested.

The formal inference and multiplicity plan, unlike the hypotheses, was
specified after the analyses had been run (Protocol Amendment 1). The six
hypotheses form one
confirmatory family controlled by Holm--Bonferroni correction. Other analyses
are classified as confirmatory, exploratory or descriptive;
Section~\ref{sec:inference} defines these classifications and the
corresponding multiplicity procedures.


\section{Contributions}
\label{sec:contrib}

To the best of our knowledge, and within the limits of the structured narrative
search described in Section~\ref{sec:search_strategy}, the study makes seven
contributions.

\begin{enumerate}
    \item \textbf{C1 (Label variation).} Three Sepsis-3 operationalisations are
    evaluated with a discrete-time competing-risks model on MIMIC-IV.
    Label-related and model-related AUROC spreads are estimated on the same test
    rows, each with an uncertainty interval, so that the two can be compared
    directly. Label sensitivity itself is established
    \cite{cohen2024subtle, dutta2026performance}; the contribution is this
    combination.

    \item \textbf{C2 (Onset timing and treatment anchoring).} The study proves
    that a conjunction-based Sepsis-3 onset rule, used here and widely in the
    MIMIC-IV literature, leaves the pre-treatment evaluation window structurally
    empty. The standard $\min(t_{\text{susp}}, t_{\text{SOFA}})$ timing rule is
    then evaluated empirically on the same cohort in a post-hoc analysis. C2
    concerns the logical consequence of the onset-timing rule.

    \item \textbf{C3 (Utility ceiling).} The adapted Utility Score is evaluated
    across the full threshold grid rather than at a single conventional
    cut-off, to test whether the low utility reported by Wang et al.\
    \cite{wang2025srpm} persists at every operating point in a controlled
    single-study comparison. Under the primary label no model exceeds
    \pcUtilityMaxGbtB{}, on a scale whose upper bound is 1.

    \item \textbf{C4 (Coefficient stability).} Cross-label coefficient
    stability is evaluated with pre-registered sign and magnitude criteria,
    testing Lauritsen et al.'s \cite{lauritsen2021framing} framing-instability
    finding on MIMIC-IV, followed by an exploratory comparison of each flagged
    change against its cluster-robust standard error.

    \item \textbf{C5 (Equity by label).} Subgroup discrimination and subgroup
    variation in who is labelled septic are measured on the same strata of an
    onset-prediction cohort under common false-discovery control. Wang et al.\
    \cite{wang2022equity} reported both quantities, but for mortality
    prediction among patients already identified as septic.

    \item \textbf{C6 (Reproducible pipeline).} A reproducible, seed-fixed
    pipeline built around a pre-registered study protocol, with every protocol
    amendment disclosed, and reported following TRIPOD+AI
    \cite{collins2024tripodai} and PROBAST+AI \cite{moons2025probastai}.

    \item \textbf{C7 (Anchor attribution).} In a post-hoc analysis, the
    Sepsis-3 conjunction is separated into its treatment-anchor and
    organ-dysfunction components, each labelled separately, to distinguish
    definitional label variation from empirical differences in membership,
    onset timing and predictability. Where C2 concerns the logic of the onset
    rule, C7 attributes the label effect empirically to its two components.
\end{enumerate}


\section{Thesis Structure}

Chapter~\ref{ch:review} reviews the evidence motivating the study.
Chapter~\ref{ch:methodology} describes the data, labels, models, validation
design and statistical analysis, and Chapter~\ref{ch:design} presents the
pipeline architecture. Chapter~\ref{ch:evaluation} reports the results,
organised around the six hypotheses; Chapter~\ref{ch:discussion} interprets
them and discusses their limitations; and Chapter~\ref{ch:conclusion}
summarises the conclusions and future work. Appendix~\ref{app:code} gives the
source-code repository and the procedure for reproducing every reported
number, and Appendix~\ref{app:stages} documents the pipeline stages.
